\documentclass{article}

\usepackage{amsmath,amssymb}
\usepackage{xurl}
\usepackage[hidelinks]{hyperref}
\setlength{\emergencystretch}{2em}

% Shared commands for notes in docs/paper-gaps/.

\DeclareUnicodeCharacter{2080}{\ifmmode{}_{0}\else\textsubscript{0}\fi}
\DeclareUnicodeCharacter{2081}{\ifmmode{}_{1}\else\textsubscript{1}\fi}
\DeclareUnicodeCharacter{2082}{\ifmmode{}_{2}\else\textsubscript{2}\fi}
\DeclareUnicodeCharacter{2099}{\ifmmode{}_{n}\else\textsubscript{n}\fi}

\newcommand{\Lean}{\textsc{Lean}}
\ExplSyntaxOn
\NewDocumentCommand{\leanid}{m}{
  \ifmmode
    \textsf{\detokenize{#1}}
  \else
    \group_begin:
    \urlstyle{sf}
    \tl_set:Nn \l_tmpa_tl {#1}
    \tl_replace_all:Nnn \l_tmpa_tl {\_} {_}
    \exp_args:NV \path \l_tmpa_tl
    \group_end:
  \fi
}
\ExplSyntaxOff
\providecommand{\tr}{\operatorname{tr}}
\newcommand{\tnleanIssueBase}{https://github.com/LionSR/TNLean/issues/}
\newcommand{\tnleanPrBase}{https://github.com/LionSR/TNLean/pull/}
\newcommand{\tnleanDocBase}{https://sirui-lu.com/TNLean/docs/}
\newcommand{\ghissue}[1]{\href{\tnleanIssueBase#1}{GitHub issue \##1}}
\newcommand{\ghpr}[1]{\href{\tnleanPrBase#1}{PR \##1}}
\newcommand{\doclink}[1]{\href{\tnleanDocBase#1.html}{\path{#1}}}

% Dirac notation and the identity operator, as in blueprint/src/macros/common.tex.
% \providecommand keeps note-local definitions authoritative.
\usepackage{dsfont}
\providecommand{\ket}[1]{|#1\rangle}
\providecommand{\bra}[1]{\langle #1|}
\providecommand{\braket}[2]{\langle #1 | #2 \rangle}
\providecommand{\ketbra}[2]{|#1\rangle\!\langle #2|}
\providecommand{\Id}{\mathds{1}}
\providecommand{\one}{\mathds{1}}
\providecommand{\C}{\mathbb{C}}
\providecommand{\MN}[1]{M_{#1}(\C)}
\providecommand{\MD}{M_D(\C)}

% External referencing for paper sources.  The build helper writes sanitized
% auxiliary files under build/paper-aux/ and excludes duplicated source labels.
\usepackage{xr}

% Import a generated paper auxiliary file with a caller-supplied label prefix.
% The candidate paths support builds from docs/paper-gaps/, the repository root,
% and build/paper-aux/ itself.
\newcommand{\importpaperaux}[2]{%
  \IfFileExists{../../build/paper-aux/#2.aux}{%
    \externaldocument[#1]{../../build/paper-aux/#2}%
  }{%
    \IfFileExists{build/paper-aux/#2.aux}{%
      \externaldocument[#1]{build/paper-aux/#2}%
    }{%
      \IfFileExists{#2.aux}{%
        \externaldocument[#1]{#2}%
      }{}%
    }%
  }%
}

% General external reference macros
\newcommand{\paperref}[2]{\ref{#1-#2}}
\newcommand{\papereqref}[2]{(\ref{#1-#2})}

% CPSV16 source (1606.00608 / MPDO-22-12-17-2.tex) external document and shortcuts
\importpaperaux{cpsv16-}{1606.00608/MPDO-22-12-17-2-tnlean}
\newcommand{\cpsvref}[1]{\paperref{cpsv16}{#1}}
\newcommand{\cpsveqref}[1]{\papereqref{cpsv16}{#1}}

% Verdict marker: \gapnote{<kind>}{<status>}. Kind is the policy
% classification (clarification, local-correction, scope-restriction,
% unfaithful, false-source, open-gap); status is open, wip (elimination
% underway), resolved, or historical. Machine-read by the site index;
% prints a verdict line.
\newcommand{\gapnote}[2]{\par\noindent\textbf{Verdict:} #1 (#2).\par}


\title{Basis-of-Representatives Scope Restriction in the\\Fundamental Theorem Formalization}
\author{}
\date{2026-05-10}

\begin{document}
\maketitle
\gapnote{scope-restriction}{resolved}

\section{The source argument}

Corollary~\texttt{II\_cor2} of
Ref.~\cite{Cirac2016MPDO_arXiv} states the general proportional Fundamental
Theorem for two tensors $A$ and $B$ whose matrix-product vectors are
proportional and which are given in basis-of-normal-tensor (BNT) canonical
form.

For $A$, let $g_A$ be the number of modulus classes, let $r_{A,j}\geq1$ be the
number of representative blocks in class $j$, let $\mu_{A,j}$ be its common
weight, and let $A_{j,q}$ denote its blocks. The source canonical form is
\begin{align}
    A^i
        &= \bigoplus_{j=1}^{g_A}
            I_{r_{A,j}}\otimes
            \mu_{A,j}
            \left(A_{j,1}^i\oplus\cdots\oplus A_{j,r_{A,j}}^i\right).
    \label{eq:cpsv16_ft_one_copy_scope_restriction_source_bnt_canonical_form}
\end{align}
The blocks in each class share the modulus $|\mu_{A,j}|$, and the classes are
strictly ordered:
\begin{align}
    |\mu_{A,1}|
        &>|\mu_{A,2}|>\cdots>|\mu_{A,g_A}|.
    \label{eq:cpsv16_ft_one_copy_scope_restriction_source_strict_modulus_order}
\end{align}
The tensor $B$ has the analogous structure, with multiplicities $r_{B,j}$ and
weights denoted by $\nu_{B,j,q}$ when individual copies are displayed.

The proof of Corollary~\texttt{II\_cor2} in
Ref.~\cite{Cirac2016MPDO_arXiv} proceeds in several steps. Step~3,
lines~1156--1163 of the arXiv version, derives for every modulus class $j$ and
every positive integer $N$ the blockwise power-sum identity
\begin{align}
    \sum_{q=1}^{r_{A,j}}\mu_{A,j,q}^N
        &= \sum_{q=1}^{r_{B,j}}
            (\nu_{B,j,q}e^{i\phi_j})^N.
    \label{eq:cpsv16_ft_one_copy_scope_restriction_blockwise_power_sum_identity}
\end{align}
The derivation uses BNT linear independence and equality of the corresponding
matrix-product vectors.

Step~4 applies Lemma~\texttt{Lem:app\_simple} of
Ref.~\cite{Cirac2016MPDO_arXiv}. For two finite sequences
$(\lambda_{a,k})_{k=1}^{x_a}$ and
$(\lambda_{b,k})_{k=1}^{x_b}$, the lemma assumes
\begin{align}
    \sum_{k=1}^{x_a}\lambda_{a,k}^N
        &= \sum_{k=1}^{x_b}\lambda_{b,k}^N
        \qquad
        (1\leq N\leq\max(x_a,x_b)).
    \label{eq:cpsv16_ft_one_copy_scope_restriction_finite_power_sum_hypothesis}
\end{align}
It concludes that the lengths agree and the sequences define the same
multiset:
\begin{align}
    x_a
        &=x_b,
    \label{eq:cpsv16_ft_one_copy_scope_restriction_power_sum_cardinality}\\
    \{\lambda_{a,k}\}_{k=1}^{x_a}
        &=\{\lambda_{b,k}\}_{k=1}^{x_b}
        \quad\text{as multisets}.
    \label{eq:cpsv16_ft_one_copy_scope_restriction_power_sum_multiset}
\end{align}
Applying this lemma to
Eq.~(\ref{eq:cpsv16_ft_one_copy_scope_restriction_blockwise_power_sum_identity})
gives
\begin{align}
    r_{A,j}
        &=r_{B,j}=:r_j,
    \label{eq:cpsv16_ft_one_copy_scope_restriction_per_class_multiplicity_match}\\
    \{\mu_{A,j,q}\}_q
        &=\{e^{i\phi_j}\nu_{B,j,q}\}_q
        \quad\text{as multisets}.
    \label{eq:cpsv16_ft_one_copy_scope_restriction_per_class_weight_match}
\end{align}
Step~6 then assembles the global gauge isomorphism as
\begin{align}
    Y
        &=\bigoplus_jI_{r_j}\otimes Y_j,
    \label{eq:cpsv16_ft_one_copy_scope_restriction_source_global_gauge}
\end{align}
using the dimension match in
Eq.~(\ref{eq:cpsv16_ft_one_copy_scope_restriction_per_class_multiplicity_match}).

The source argument therefore recovers the multiplicities and weight multisets
from finitely many polynomial identities. It does not use a limiting argument.

\section{The basis-of-representatives restriction}

The surviving formal normal-CF-BNT structure is
\leanid{IsNormalCanonicalFormBNT}.\footnote{Defined in
\path{TNLean/MPS/BNT/Construction.lean}. The earlier non-normal wrapper was
retired after its downstream consumers were replaced by explicit
separated-block hypotheses.} This structure no longer requires strict
decrease of the weight moduli. It requires only non-increasing order by
modulus, retains only nonzero weights, and permits equal-modulus blocks.
Instead of strict modulus separation, its BNT separation field requires
distinct blocks not to be gauge-phase equivalent.

The remaining restriction is therefore not strict modulus ordering.
The structure \leanid{IsNormalCanonicalFormBNT} is a
basis-of-representatives surface: its block family already contains one
selected representative for each gauge-phase class. It does not retain
multiple copies of one source BNT tensor as separate summands together with
their individual weights $\mu_{j,q}$, copy gauges, and coefficient
\begin{align}
    c_N^{(j)}
        &=\sum_q\mu_{j,q}^N.
    \label{eq:cpsv16_ft_one_copy_scope_restriction_sector_copy_coefficient}
\end{align}
The formalization carries this raw multiplicity data in
\leanid{SectorDecomposition} and in the SectorBNT comparison theorems, not in
\leanid{IsNormalCanonicalFormBNT}.

In the one-copy special case, the source decomposition becomes
\begin{align}
    A^i
        &=\bigoplus_{j=1}^{g}\mu_jA_j^i.
    \label{eq:cpsv16_ft_one_copy_scope_restriction_one_copy_decomposition}
\end{align}
The power-sum identity then reduces to
\begin{align}
    \mu_{A,j}^N
        &=(e^{i\phi_j}\nu_{B,j})^N.
    \label{eq:cpsv16_ft_one_copy_scope_restriction_one_copy_power_identity}
\end{align}
Steps~3--4 of the source proof become vacuous, and the global gauge in Step~6
reduces to
\begin{align}
    Y
        &=\bigoplus_jY_j.
    \label{eq:cpsv16_ft_one_copy_scope_restriction_one_copy_global_gauge}
\end{align}
The present normal-CF-BNT surface is a representative-family surface of this
kind. The full copy-multiplicity comparison is supplied elsewhere.

\section{Components absent from the representative surface}

Four components are absent from the normal-CF-BNT representative surface.
Each is required to recover the source-level multiplicity data with
$r_j\geq1$.

\subsection{Finite Newton--Girard uniqueness}

Lemma~\texttt{Lem:app\_simple} of Ref.~\cite{Cirac2016MPDO_arXiv} states that
two finite sequences whose power sums agree for
$N=1,\ldots,\max(x_a,x_b)$ are equal as multisets. This finite algebraic fact
follows from the Newton--Girard identities: the power sums determine the
elementary symmetric polynomials, which determine the multiset of roots.

A module named \leanid{ScalarPowerSumIdentity}\footnote{Introduced in commit
\texttt{dd54b79e}.} was added early in the project, but it does not specialize
to the fixed-length statement in the form required by Step~4 of the source
proof. The result must therefore be available as a standalone theorem,
separate from the BNT context.

\subsection{BNT power-sum identity}

The identity in
Eq.~(\ref{eq:cpsv16_ft_one_copy_scope_restriction_blockwise_power_sum_identity})
requires a derivation from the BNT structures on $A$ and $B$ and the equal-MPV
hypothesis. For every modulus class $j$, equality of the MPVs and BNT linear
independence must imply equality of the corresponding power sums.

The formal development no longer packages the coefficient arrays of the
discarded proportional branch as standalone hypotheses. On the representative
surface, the source derivation of
Eq.~(\ref{eq:cpsv16_ft_one_copy_scope_restriction_blockwise_power_sum_identity})
is therefore not carried out. That derivation is logically prior to the use
of Lemma~\texttt{Lem:app\_simple} of
Ref.~\cite{Cirac2016MPDO_arXiv}.

\subsection{Multi-copy sector structure}

The structure \leanid{IsNormalCanonicalFormBNT} permits equal moduli but does
not store repeated gauge-phase-equivalent copies with their individual
weights. The general theorem therefore requires the
\leanid{SectorDecomposition}/SectorBNT multiplicity structure, in which every
sector has an arbitrary finite list of copy weights and gauges. Without this
data, the general-multiplicity form of Corollary~\texttt{II\_cor2} of
Ref.~\cite{Cirac2016MPDO_arXiv} and its gauge-assembly step cannot be stated on
the normal-CF-BNT surface alone.

\subsection{Global gauge assembly}

For $r_j\geq1$ copies per sector, the global gauge has the form in
Eq.~(\ref{eq:cpsv16_ft_one_copy_scope_restriction_source_global_gauge}).
Constructing and verifying this gauge requires the dimension equality in
Eq.~(\ref{eq:cpsv16_ft_one_copy_scope_restriction_per_class_multiplicity_match}),
which depends on both finite power-sum uniqueness and the BNT power-sum
identity. The simpler gauge in
Eq.~(\ref{eq:cpsv16_ft_one_copy_scope_restriction_one_copy_global_gauge}) is
already present.

\section{What the current formalization proves}

On the basis-of-representatives surface, the formalization establishes the
equal-MPV branch of the Fundamental Theorem argument. BNT separation gives
cross-block overlap decay, and an equal-MPV hypothesis at one witness length
matches the per-block gauges.

The former strict-modulus proportional branch has been deleted. The surviving
block-matching boundary uses the explicit conclusion
\leanid{BlockPermutationGaugePhaseConclusion} only when the BNT comparison has
already been supplied. Thus no current theorem is presented as deriving the
source proportional comparison from the discarded coefficient hypotheses.

Neither the old dominant-normalization field nor the old strict-ordering field
remains in \leanid{IsNormalCanonicalFormBNT}. When a CPSV argument needs a
unit-modulus witness for a sector, it carries that datum as an explicit
SectorBNT hypothesis, for example through the corresponding field of
\leanid{SectorDecomposition}.

\section{Downstream audit status}

The downstream CPSV16 surfaces that consume the restricted normal-CF-BNT
predicate have been audited against the basis-of-representatives
boundary.\footnote{The relevant \Lean{} surfaces are
\leanid{IsNormalCanonicalFormBNT},
\leanid{IsNormalCanonicalFormBNT.ofSeparatedData},
\leanid{IsNormalCanonicalFormBNT.cross_overlap_tendsto_zero},
\leanid{IsNormalCanonicalFormBNT.isBNT},
\leanid{IsNormalCanonicalFormBNT.exists_common_blockTensor_isInjective},
and
\leanid{exists_common_blockTensor_isInjective_two_of_isNormalCanonicalFormBNT}.}

Each public surface assumes or constructs an already selected representative
family whose weights are ordered non-increasingly by modulus. Equal moduli are
allowed, but this surface does not reconstruct repeated
gauge-phase-equivalent copies or their individual weights. These statements
neither introduce a source-faithful canonical-form existence theorem nor
recover the repeated-copy sector data in the BNT decomposition of
Ref.~\cite{Cirac2016MPDO_arXiv}. Their docstrings and the corresponding
Wielandt blueprint corollary therefore carry the same
basis-of-representatives marker.

This marker does not apply to the sector-BNT development based on
\leanid{IsBNTCanonicalForm}. That predicate is the paper-faithful sector
surface used for the CPSV16 coefficient and weight-multiplicity arguments,
where repeated copies inside one sector remain explicit.

\section{Conclusion}

The representative-family formalization proves the proportional Fundamental
Theorem on an already selected basis-of-representatives surface. The general
case of Corollary~\texttt{II\_cor2} in
Ref.~\cite{Cirac2016MPDO_arXiv}, including repeated copies within a sector,
has since been supplied on the SectorBNT surface.

The standalone Newton--Girard uniqueness result is
\leanid{Matrix.sum_pow_eq_implies_card_eq_and_multiset_eq_of_le_max_card}.
The BNT coefficient identity derived from the equal-MPV hypothesis is
\leanid{MPSTensor.coeff_identity_via_global_gauge}, using the phase-valued
auxiliary theorem
\leanid{MPSTensor.coeff_identity_via_matched_mpv_phase}. The power-sum
expression and weight recovery are supplied by
\leanid{MPSTensor.SectorDecomposition.coeff_eq_sum_weight_pow},
\leanid{MPSTensor.matched_sector_weight_multiset_eq}, and
\leanid{MPSTensor.matched_sector_weight_equiv}. The global gauge with identity
factors on the multiplicity spaces is provided by
\leanid{MPSTensor.ft_sector_bnt_equal_global_gauge}.

The corresponding tracked items,
\href{https://github.com/LionSR/TNLean/issues/1560}{issue~\#1560},
\href{https://github.com/LionSR/TNLean/issues/1561}{issue~\#1561}, and
\href{https://github.com/LionSR/TNLean/issues/1562}{issue~\#1562}, are closed.

This note therefore records why the older one-copy surface remains restricted.
The paper-faithful multi-copy route is the
\leanid{MPSTensor.IsBNTCanonicalForm}/SectorBNT declaration path and does not
carry the one-copy scope-restriction marker. The local scope-restriction marker
on the one-copy BNT comparison in the blueprint remains correct.

\bibliographystyle{alpha}
\bibliography{references}

\end{document}
